\documentclass[10pt,a4paper]{amsart}
\usepackage[margin=19mm]{geometry}
\usepackage{amsmath,amssymb,mathtools}
\usepackage[scr=rsfs,cal=euler]{mathalfa}
\usepackage{xfrac}
\usepackage[unicode,hidelinks,bookmarks=false]{hyperref}
\newcommand{\ie}{i.e.\ }
\newcommand{\cf}{cf.\ }
\newcommand{\resp}{respectively\ }
\newcommand{\Subsection}{\subsection}
\providecommand{\nobreakdash}{\nobreak\hbox{-}}
\setlength{\emergencystretch}{2em}
\multlinegap0pt
\usepackage{amssymb}
\usepackage{mathtools}
\usepackage[normalem]{ulem}
\newtheorem{lemma}{Lemma}
\newcommand{\R}{\mathbb{R}}
\newcommand{\Z}{\mathbb{Z}}
\newcommand{\e}{\varepsilon}
\newcommand{\ez}{\e\Z^N}
\newcommand{\g}{\gamma}
\newcommand{\sd}{\,\mathrm{sd}}
\newcommand{\ud}{\,\mathrm{d}}
\title{Elementary discrete diffusion/redistancing schemes for the mean curvature flow}
\author{Antonin Chambolle, Daniele De Gennaro, Massimiliano Morini}
\date{Source: arXiv:2506.05946v2 (CC BY 4.0)}
\begin{document}
\maketitle

\noindent\emph{Source and licence.} Elementary discrete diffusion/redistancing schemes for the mean curvature flow. Version arXiv:2506.05946v2 (CC BY 4.0), distributed by arXiv under the Creative Commons Attribution 4.0 International licence, http://creativecommons.org/licenses/by/4.0/, as recorded in the arXiv licence metadata for this exact version and shown on the abstract page https://arxiv.org/abs/2506.05946v2. Full licence notice: \texttt{licenses/arxiv-2506.05946-CC-BY-4.0.txt}.\par\smallskip

\noindent\textit{Editorial scope.} Counted unit: the article's lemma lem:evol\_balls\_g (source0.tex lines 1176--1181). The article's complete original proof of it is delivered with it (source0.tex lines 1183--1272, one \texttt{proof} environment of 90 lines). No mathematics is written, repaired or re-derived: the statement and the proof are the article's own lines, in the article's own notation, and no retained line is substituted or re-flowed.  Retained uncounted as the source-local context the proof needs: lines 211--216; lines 236--248; lines 459--525; lines 482--487; lines 765--769.  Nothing was omitted from any retained span, so the package carries no omission record.  No retained line contains a citation, so the wrapper carries no bibliography and no bibliography entry.  No retained line contains a figure environment, an image-inclusion command or drawing code, so no image is delivered and the package declares no asset dependency.  Editorial formatting only, in addition: an AMS \texttt{amsart} preamble that restates the article's own declarations and macros used by the retained lines, namely the environments \texttt{lemma} on separate counters, together with the standard packages those lines need.  The retained environments print in this document as Lemma 1 in order of appearance, whereas the article prints the same items with the numbering of its own full document.  The complete native source of the article as distributed in the arXiv e-print for this exact version is bundled unchanged as \texttt{sources/179426/source0.tex}.
\begin{equation}\label{eq:sdp}
  \begin{cases}
    \ud^+[u]_i := \inf_{j:u_j<0}  {(} u_j + |j-i| {)}\,,\\
    \sd^+[u]_i := \sup_{j:u_j\ge 0}  {(} \ud^+_j - |j-i|  {)}
  \end{cases}
\end{equation}

\begin{lemma}\label{lem:defsd}
 {Assume still that $(u_i)_{i\in\ez}$ is $1$-Lipschitz. Then}
  the distance function $\sd^+$ may also be defined as follows:
  \begin{enumerate} 
  \item $\ud^+$ is the largest $1$-Lipschitz function such that $\ud^+_i\le u_i$ when
    $u_i<0$, and in particular $\ud^+_i=u_i$ if $u_i<0$ and $\ud^+_i\ge u_i$ if $u_i\ge 0$,
  \item $\sd^+$ is the smallest $1$-Lipschitz function such that $\sd^+_i\ge \ud^+_i$
    when $u_i\ge 0$, and in particular $\sd^+_i=\ud^+_i\ge u_i$ if $u_i\ge 0$ and $\sd^+_i\le \ud^+_i=u_i$
    if $u_i<0$.
  \end{enumerate}
  We also deduce that $\sd_i^+\ge 0\Leftrightarrow u_i\ge 0$ and
  $\sd_i^+< 0 \Leftrightarrow u_i<0$.
\end{lemma}

As said, a crucial point for proving the convergence of the method, is to control the behavior
of the algorithm when $\sd^{k,\e}_i$ represents a ball of radius $R>0$. For this,
given $\theta\in  {(0,1]}$, we let for $i\in\ez$,
\begin{equation}\label{eq:v}
    u_i= |i|-R,\quad v_i = (1-\theta) u_i + \frac{\theta}{2N} \sum_{n=1}^N (u_{i+\e e_n}+u_{i-\e e_n})   
\end{equation}
and let $d=\sd^+[v]$. Then, we show:
\begin{lemma}\label{lem:balls}
  There exists a dimensional constant $C\ge 1$ such that if $\e/R$ is small enough (depending
  only on the dimension), then for every $i\in\ez$
  \[
    d_i \le |i|-R + \frac{C}{R}\e^2 = |i|-R + \frac{2NC}{\theta R}h  
  \]
\end{lemma}
\noindent Here as before $h,\e,\theta$ are linked by the relationship $\theta=2Nh/\e^2$.
\begin{proof}
  Without loss of generality we assume $\theta=1$. 
  We first observe that
  \[
    v_i = -R + \frac{1}{2N}\sum_{n=1}^N |i+\e e_n|+ |i-\e e_n|.
  \]
  We remark that if $|i|\ge 2\e$, $e\in\{e_n,-e_n:n= 1,\dots ,N\}$,
   {by a   Taylor expansion it holds:}
  \begin{multline}\label{eq:Taylor}
    |i+\e e| = |i| + \e\frac{i\cdot e}{|i|} +
    \int_0^1 (1-t) \frac{1}{|i+t\e e|}\left(I - \frac{(i+t\e e)\otimes (i+t\e e)}{|i+t\e e|^2}
    \right)\e e \cdot \e e\,dt  \\ 
    \le  |i| + \e\frac{i\cdot e}{|i|} + \frac{1}{|i|-\e}\frac{\e^2}{2}
  \end{multline}
  so that:
  \[
    v_i \le |i|-R + \frac{1}{|i|-\e}\frac{\e^2}{2}, \qquad  {\text{for } |i|\ge 2\e}.
  \]
  If we assume that $\e\le \min\{1,R/2\}$, then for $|i|\le R-\e$,
  $v_i< 0$. Hence, since by Lemma~\ref{lem:defsd}
  it is enough
  to estimate $d_i^+$ at points
  $i\in\ez$ with $v_i\ge 0$, we assume  $|i|\ge R-\e$. For all such $i$ we have:
  \begin{equation}
  \begin{split}\label{eq:d+est}
      \ud^+[v]_i :&= \inf_{j:v_j<0}  {(} v_j + |j-i|  {)}
    \le \inf_{j: \frac{R}{2}+\e\le |j|\le R-\e}  {(} v_j + |j-i|  {)} \\
    &\le \inf_{j: \frac{R}{2}+\e \le |j|\le R-\e}  {(} |j|-R + \frac{\e^2}{R} + |j-i|  ),
  \end{split}
  \end{equation}
  assuming $\e\le R/8$ so that the set of $j$'s is not empty.
  Consider $j$ with $\frac{R}{2}+\e \le |j|\le R-\e$, close to the
  segment $[0,i]$:  if $\tilde{\jmath}$ is the projection of $j$ onto $[0,i]$, one has
  \begin{multline*}
    |j|+|j-i| =  \sqrt{ |j-\tilde{\jmath}|^2 +|\tilde{\jmath}|^2} + \sqrt{|\tilde{\jmath}-i|^2+|j-\tilde{\jmath}|^2}
    \\
    \le |\tilde{\jmath}| + \frac{|j-\tilde{\jmath}|^2 }{2|\tilde{\jmath}|}
    + |\tilde{\jmath}-i| + \frac{|j-\tilde{\jmath}|^2 }{2|\tilde{\jmath}-i|}
    = |i| + \left(\frac{1}{2|\tilde{\jmath}|} + \frac{1}{2|\tilde{\jmath}-i|} \right){|j-\tilde{\jmath}|^2 }
  \end{multline*}
  If $\e/R$ is small enough (depending only on $N$), we can find
  $j,\tilde{\jmath}$ such that $|j-\tilde{\jmath}|^2\le N\e^2$ and $R/2\le |\tilde{\jmath}|\le 3R/4$,
  so that $|\tilde{\jmath}-i|\ge R/4$. We obtain for such a choice:
\begin{equation}\label{eq:triangle}
    |j|+|j-i| \le |i|+\frac{3N}{R}\e^2
\end{equation}
  This shows that where $d^+_i$ is non-negative, it is
  less than $|i|-R + \frac{3N+1}{R}\e^2$ as soon as $\e/R$ is small enough. As $\sd^+$ is the smallest
  $1$-Lipschitz function larger than $d^+_i$ where it is non-negative (Lemma~\ref{lem:defsd}),
  this achieves the proof.
\end{proof}
\subsection{Consistency of the algorithm}\label{sec:consistexplicit}

  \begin{multline}
    |i+\e e| = |i| + \e\frac{i\cdot e}{|i|} +
    \int_0^1 (1-t) \frac{1}{|i+t\e e|}\left(I - \frac{(i+t\e e)\otimes (i+t\e e)}{|i+t\e e|^2}
    \right)\e e \cdot \e e\,dt  \\ 
    \le  |i| + \e\frac{i\cdot e}{|i|} + \frac{1}{|i|-\e}\frac{\e^2}{2}
  \end{multline}

\begin{align}
  & \textstyle  K_j^\e= K_{-j}^\e \ge 0 \quad \text{ for all } j\in\e\Z^N, \label{Kpositive} \\
   & \textstyle \sum_{j\in\e\Z^N} K_j^\e=1, \label{Kone}\\
  & \textstyle h=h(\e):= \frac{1}{2N}\sum_{j\in\e\Z^N} |j|^2K_j^\e \to 0 \quad \text{  as } \quad  \e\to 0.  \label{Kh}  
\end{align}

\begin{lemma}\label{lem:evol_balls_g}
    Assume that $h(\e)\ge \e^2$.  There exist two constants $R_0,c>0$, such that the following holds. For every $R\in (0,R_0)$ and  for $h$ small enough  it holds
\begin{equation}\label{eq:speed_ball_g}
\sd^+\Big[\gamma^{-1}(K^\e\ast \gamma(|\cdot|-R))\Big]\le |\cdot|-R + c\tfrac hR.
\end{equation}
\end{lemma}

\begin{proof}
Let us fix $0<R_0\le 1$ such that $\g'(R_0)\ge \frac12$. We denote 
\[
v:=K^\e\ast \g(|\cdot|-R).
\]
We start remarking that  the thesis 
follows once we show that in the region $\{ v\ge 0\}$ it holds 
\begin{equation}\label{eq:goal_sd+_g}
    \text{d}^+[\g^{-1}( v)]\le |\cdot|-R+\tfrac{ch}R,
\end{equation}
as then $\text{sd}^+$ is the smallest 1-Lipschitz function lying above $\text{d}^+$ where it is non-negative. 



We will in fact prove the following inequality
\begin{equation}\label{eq:subgoal_sd+_g}
    v\le \g(|\cdot|-R+\tfrac{ch}R)  \qquad \text{for all } |i|\in (\tfrac R4, \tfrac R2 ),
\end{equation}
for $R\le R_0$ and $h,\e$ small. 


 
\noindent\textit{Proof of \eqref{eq:speed_ball_g} assuming \eqref{eq:subgoal_sd+_g}.}
If \eqref{eq:subgoal_sd+_g} holds, then we deduce that  
\begin{equation}\label{eq:sign_v}
    v\le 0\quad  \text{ in } B_{R/2}\cap \ez.
\end{equation}
Indeed, from \eqref{eq:subgoal_sd+_g} this inequality already holds in the region $(B_{R/2}\setminus B_{R/4})\cap \ez$. Then, consider a point $i \in B_{R/4} \cap \ez$: there exists $j$ such that $|i - j| \le R/4 + \sqrt{N}\, \e$ and $  |j| \in (R/4,R/4 + \sqrt{N}\, \e)$. By the 1-Lipschitz property, we obtain
\[
v_i \le v_j + \tfrac{R}{4} + \sqrt{N}\, \e \le \g\left(-\tfrac{3R}{4} + \sqrt{N}\, \e + \tfrac{ch}{R}\right) + \tfrac{R}{4} + \sqrt{N}\, \e.
\]
Moreover, assuming $\e,h$  small so that $\g'(-\tfrac{3R}{4} + \sqrt{N}\, \e + \tfrac{ch}{R}) \ge \g'(-R_0) \ge 1/2$, by the equation above and the  convexity of $\g$ in $\{\g \le 0\}$ we get 
\[
\begin{split}
v_i \le \g\left(-\tfrac{3R}{4} + \sqrt{N}\, \e + \tfrac{ch}{R}\right) + \g'\left(-\tfrac{3R}{4} + \sqrt{N}\, \e + \tfrac{ch}{R}\right)\left(\tfrac{R}{2} + 2\sqrt{N}\, \e\right)\\
\le \g\left(-\tfrac{R}{4} + 3\sqrt{N}\, \e + \tfrac{ch}{R}\right).
\end{split}
\]
The right-hand side can be made negative by choosing $\e, h$ sufficiently small. This concludes the proof of \eqref{eq:sign_v}.

With \eqref{eq:sign_v} in hand, we can  follow the argument from Lemma~\ref{lem:balls} to establish \eqref{eq:goal_sd+_g}, and thus complete the proof of the lemma. Specifically, since \eqref{eq:sign_v} holds in $B_{R/2}\cap \ez$, it remains to prove \eqref{eq:goal_sd+_g} for $|i| \ge R/2$. For such $i$, we apply the definition of $\text{d}^+$ from \eqref{eq:sdp}, restricting the infimum to the region where $|j| \in (R/4, R/2)$. Using the bound from \eqref{eq:subgoal_sd+_g}, we obtain an estimate analogous to the right-hand side of \eqref{eq:d+est}. Since $|i|$ is bounded from below, the conclusion follows exactly as in Lemma~\ref{lem:balls}.



\noindent\textit{Proof of \eqref{eq:subgoal_sd+_g}.} To show \eqref{eq:subgoal_sd+_g}, we fix $i\in\ez$ with $|i|\in (R/4,R/2)$  and  estimate 
\begin{equation}
    \begin{split}\label{eq:est_v_g}
    v_i=(K^\e\ast \g(|\cdot|-R))_i &\le  \sum_{|j|\le R/8} K^\e_j \g(|i-j|-R) +  \sum_{|j|\ge R/8} K^\e_j  \Big( \g(|i|-R)+|j|\Big) \\
    &\le \sum_{|j|\le R/8} K^\e_j \g(|i-j|-R) +  \sum_{|j|\ge R/8} K^\e_j \g(|i|-R) +\tfrac {8h}{R},
    \end{split}
\end{equation}
where  in the first inequality we used that $\g$ is 1-Lipschitz and in the second one we used \eqref{Kh}. We need to estimate the term
\begin{equation}\label{eq:first_term}
    \sum_{|j|\le R/8} K^\e_j \g(|i-j|-R).
\end{equation}
Recalling \eqref{eq:Taylor}, for $i,j\in\ez$ with  $|i|\ge R/4$ and $|j|\le R/8$,   {by a Taylor expansion} it holds 
\[
    |i\pm j|\le |i| \pm \frac{i\cdot j}{|i|} + \frac1{|i|-|j|}\frac{|j|^2}2 \le |i| \pm \frac{i\cdot j}{|i|} + \frac 4R |j|^2.
\] 
Therefore, we  can rewrite the term \eqref{eq:first_term} as follows:
\begin{equation}
    \begin{split}
        \sum_{|j|\le R/8} K^\e_j \g(|i-j|&-R) =\frac 12  \sum_{|j|\le R/8} K^\e_j (\g(|i-j|-R) +\g(|i+j|-R) ) \\
    &\le \frac 12  \sum_{|j|\le R/8} K^\e_j (\g(|i|-R -\tfrac{j\cdot i}{|i|} + \tfrac{4}{R} |j|^2 ) +\g(|i|-R+\tfrac{j\cdot i}{|i|} + \tfrac{4}{R}|j|^2) ).\label{eq:mult_g}
    \end{split}
\end{equation}
Let $x,y\in \R$. Since $\g''$ is bounded from above, by a Taylor development we get 
\[
    \g( x+y)  = \g(x) +\g'(x)\, y + \frac{|y|^2}2 \int_0^1\g''(x +  sy)(1-s)\,ds\le \g(x) +\g'(x)\, y + c|y|^2.
\]
We  thus use this estimate to get  
\begin{multline*}
    \g(|i|-R + \tfrac{4}{R} |j|^2 \pm \tfrac{j\cdot i}{|i|}  ) \le  \g(|i|-R+ \tfrac{4}{R} |j|^2) \pm \frac{j\cdot i}{|i|}\, \g'(|i|-R+ \tfrac{4}{R} |j|^2) + c |j|^2\\
    \le    \g(|i|-R  ) \pm \frac{j\cdot i}{|i|} \, \g'(|i|-R+ \tfrac{4}{R} |j|^2) + |j|^2 (c+\tfrac 4{R})
\end{multline*}
where in the last inequality we used the 1-Lipschitz property. Plugging the inequality above in \eqref{eq:mult_g} and using that   $|i|\ge R/2$ and $R\le 1$, we have bounded the term \eqref{eq:first_term} as follows 
\[
\sum_{|j|\le R/8} K^\e_j \g(|i-j|-R)\le \sum_{|j|\le R/8} K^\e_j \g(|i|-R) + \frac{ch}R,
\]
where $c$ is a positive constant.  
Therefore, by \eqref{eq:est_v_g} we get 
\[
(K^\e\ast v)_i \le \g(|i|-R) +\frac{ch}R,
\]
where $c$ denotes a positive constant.  Since $|i|\in(R/4,R/2)$ then $\g'(|i|-R)\ge \g'(-R_0)\ge 1/2$, and therefore 
\[
(K^\e\ast v)_i \le \g(|i|-R) + 2c \frac{h}R\g'(|i|-R)\le \g(|i|-R+ \tfrac{2ch}R),
\]
where the last inequality holds by convexity as long as  $  |i|-R+\tfrac{2ch}R\le 0$, which holds for $h$ small.
\end{proof}

\end{document}
